\documentclass[tikz,svgnames,border={10 10},convert={density=300, outext=.png}]{standalone}

% 若需要中文，请用 xelatex 编译并取消下面两行注释：
\usepackage{xeCJK}
% \setCJKmainfont{Noto Sans CJK SC}

\usetikzlibrary{positioning,arrows,fit}

\begin{document}
\begin{tikzpicture}[
    box/.style={rectangle,draw,fill=DarkGray!20,rounded corners,
                node distance=1cm,text width=16em,text centered,
                minimum height=2.2em,thick},
    arrow/.style={draw,-latex',thick},
  ]

  % === 主体流程节点（与示例同样的布局与风格） ===
  \node[box,fill=PeachPuff!60,text width=20em] (input)
    {接收上游输入：事件/消息（含时间戳/优先级）};

  \node[box,below=0.6 of input] (sched)
    {事件调度（按时间戳 / 维护 GVT）};

  \node[box,below=0.5 of sched] (bus)
    {消息路由（发布 / 订阅）};

  \node[box,below=0.5 of bus] (sync)
    {并行与时间同步\\（时间窗 / 回溯保障）};

  \node[box,below=0.5 of sync] (crit)
    {仿真结束条件满足？};

  % 底部输出
  \node[box,below=1.8 of crit,fill=LightSteelBlue!60,text width=22em] (output)
    {输出：执行结果/回执、撮合/成交消息、监控指标};

  % === 箭头（直连 + 与示例相同的 Yes/No 回路） ===
  \path[arrow] (input) -- (sched);
  \path[arrow] (sched) -- (bus);
  \path[arrow] (bus)   -- (sync);
  \path[arrow] (sync)  -- (crit);

  % Yes 分支：直达输出
  \path[arrow] (crit) -- (output)
    node[midway,left=0.1,draw,outer sep=0pt,fill=white] {Yes};

  % No 分支：先竖直向下，再水平右移（No 在水平段下方），最后折返到调度节点
  \path[draw,thick] (crit.south) ++(0,-1em) coordinate(novert);
  \draw[thick] (novert) -- ++(4,0) coordinate(nojump)
        node[midway,below=2pt,draw,outer sep=0pt,fill=white] {No};
  \draw[arrow] (nojump) |- (sched.east);

\end{tikzpicture}

\end{document}